\clearpage
\section{Effective-Address Profiles}\label{page:effective-addressing-modes}

EA is the base seven-bit compact effective-address profile selected by the consuming instruction\textquotesingle{}s operand
type. Extensions may define additional profiles that reuse the memory address calculation rules and EXT1/EXT2
descriptor grammar in this section while assigning profile-specific compact forms. An effective-address (:term:base.terms.field:) is an operand
descriptor with an instruction-defined role and interpretation width. A
value role reads or writes the selected register, immediate, or memory value. An address role uses a register or
immediate as an address value and uses a memory-form EA as the calculated address without an initial data read. A
control-target role uses a register or immediate value directly and reads an interpretation-width target from a
memory form. Memory forms continue into (:term:base.terms.segment_pre_translation:) and paging.

Each compact (:term:base.terms.field:) is embedded in the instruction\textquotesingle{}s (:term:base.terms.opcode_space:), and all compact (:term:base.terms.field|plural:) precede any appended
effective-address (:term:base.terms.payload:) bytes. Any displacement, absolute address, immediate, or extended
descriptor bytes are appended as separate (:term:base.terms.payload|plural:). All extended descriptor bytes appear first, ordered according to the
declared instruction operand order of the operands that select them. All remaining effective-address (:term:base.terms.payload|plural:) follow, also in
declared instruction operand order.
For two payload-bearing operands, the encoded stream contains the (:term:base.terms.opcode_space:) for both operands, any extended descriptor for
operand 0, any extended descriptor for operand 1, the remaining operand 0 (:term:base.terms.payload|plural:), the remaining operand 1
(:term:base.terms.payload|plural:), and any trailing (:term:base.terms.padding:), in that order.

Multi-byte (:term:base.terms.payload|plural:), such as displacements, absolute addresses, and immediates, are little-endian.
The payload type declaration determines signedness: address payloads declare \texttt{signed}, and integer immediates declare
\texttt{signed_integer} or \texttt{unsigned_integer}. Signed payloads are sign-extended before use. The consuming instruction
determines any further interpretation of the value. Extended descriptor bytes appear in stream order, byte 0 first. Each
descriptor byte is independently numbered from bit 7 through bit 0.

Indexed scale is the EA interpretation width declared by the consuming instruction encoding. B, W, L, and Q widths select scales 1,
2, 4, and 8. An encoding without a fixed width uses its selected instruction size. LEA and SEGLEA use
their suffix; size-less instructions with address or control-target roles use Q.

For every compact or extended-descriptor PC-based EA, the PC term is the architectural PC naming byte 0 of the current instruction.
The PC term remains that instruction-start value throughout operand evaluation.

\clearpage
